\documentclass[11pt,a4paper]{article}
\usepackage[T1]{fontenc}
\usepackage[utf8]{inputenc}
\usepackage{times}
\usepackage[top=2cm,bottom=2cm,left=2.2cm,right=2.2cm]{geometry}
\usepackage{microtype}
\usepackage[hidelinks]{hyperref}
\pagestyle{empty}
\setlength{\parindent}{0pt}
\setlength{\parskip}{7pt}

\begin{document}

{\large\bfseries Grover Aruquipa}\par\vspace{-5pt}
{\small Qu\'ebec, QC, Canada \ \textbullet\ \ grover.aruquipa.9@gmail.com \ \textbullet\ \ groveraruquipa.github.io}

\vspace{4pt}
\textbf{Supporting statement: Postdoctoral Research Assistant in Tensegrity Robot Design and Control (vacancy 189502)}\\
Protean Robotics Lab, Department of Engineering Science, University of Oxford

Dear Professor Baines,

Thank you again for your advice on the Branco Weiss Fellowship in September; your point about societal impact changed how I think about my own projects. I am now writing to apply for the postdoctoral position in tensegrity robot design and control. My whole PhD has been about one idea: a robot can do more when the mechanism itself does part of the work, so that grasping, rotation and redundancy come from the structure rather than from devices added at the tip. Tensegrities take that idea further than any architecture I have worked with, because the structure and the robot are the same thing. Working on them in your lab is the next step I would most like to take.

\textbf{PhD in engineering.} I am completing a PhD in Mechanical Engineering at Universit\'e Laval with Prof.\ Cl\'ement Gosselin and Prof.\ Patrice Lambert, on kinematically redundant parallel robots with unlimited rotations and integrated grasping. I submit the thesis in December 2026 and defend in January 2027. I am first author of three journal articles (ASME Journal of Mechanisms and Robotics, Mechanism and Machine Theory), including a review of parallel robots for grasping, and of four further papers under review, one of them at ICRA 2027.

\textbf{Constructing and programming robots, including motor control.} I have designed, fabricated and controlled three parallel robots with 4, 7 and 8 actuators, all mounted on the base, with the gripper driven through the mechanism. For each one I did the kinematic model, the SolidWorks design, the machining, SLS and FDM printing, the wiring and the commissioning, and I wrote the multi-motor control running at 500\,Hz with incremental encoders (OPAL-RT, Simulink Real-Time, and C++ in real time during my master's research at FEMTO-ST). The most compact of them, a 4-DoF coaxial wrist-gripper with 0.25\,kg of moving mass, turned its tool through ten continuous revolutions with sub-degree tracking while holding the object. I am now building a four-leg Delta robot with an integrated gripper that estimates grasp force from motor currents, and I validated its dynamic model in Simscape before moving to hardware.

\textbf{Simulation.} I use MuJoCo and Simscape to test my mechanisms before building them, and I compare the simulated dynamics with the measured behaviour of the prototypes. I have also experimented with Isaac Sim, which gives me a starting point for IsaacLab and its learning workflows. The parts that are hard to learn, getting the mechanics of the model right and checking it against hardware, are the parts I already do every day.

\textbf{Co-optimising design and control of nonlinear structures.} In redundant parallel robots, design and control cannot be separated: the geometry decides where the singularities are, and the redundancy law decides whether the robot can avoid them. My work on collision-free and singularity-free unlimited rotation is a joint problem of this kind. The statics I use are also close to those of tensegrities. A Type II singularity of a parallel robot and an infinitesimal mechanism of a tensegrity are both rank losses of the same kind of equilibrium matrix, and a state of self-stress is a vector in the null space of that matrix. I would bring this way of reading the structure to the co-optimisation, so that analytical measures of stiffness, prestress and distance to singular shapes can guide the search over designs and policies in simulation.

\textbf{Supervising students.} I supervise two graduate students day to day on an 8-actuator (6+2)-DoF robot, from concept to a working prototype and a co-authored ICRA 2027 submission. I led the Universit\'e Laval team that won first place in the graduate division of the 2025 ASME Student Mechanisms and Robotics Design Competition, and I regularly lecture on parallel robots at universities in Bolivia.

\textbf{Experience with tensegrity systems.} I have not built a tensegrity robot. The closest work I have done is on cable-driven parallel mechanisms, where I have studied statics and force closure, and on reconfigurable platforms whose shape changes the way forces are distributed. I would come to tensegrities with the mechanics already in place and a great deal to learn about building them, which I would welcome.

\textbf{Research interests and fit.} Your lab's work on adaptive morphology, from multi-environment robots that change shape to Spikebot and the steerable jumping tensegrity, shows robots whose body is a large part of their intelligence. That is the direction I want my research to take. Two questions I would be glad to pursue within the project are, first, whether the shape change of a tensegrity can be used for manipulation as well as locomotion, grasping and reorienting objects with the structure itself, which would connect directly to my work on integrated grasping and unlimited rotation; and second, how far simple structural measures can reduce the cost of searching jointly over tensegrity designs and control policies. In the longer term I aim for an academic career, and a post in your group, with the Oxford Robotics Institute around it, is the environment where I would most like to build it. If you think it would be a good fit, I would also be glad to prepare a Branco Weiss application with your lab as host in a coming round, building on the project.

I can start in early 2027, after my defense, and I would need a UK Skilled Worker visa. My referees for this application are Prof.\ Cl\'ement Gosselin (Universit\'e Laval, clement.gosselin@gmc.ulaval.ca) and Univ.-Prof.\ Andreas M\"uller (JKU Linz, a.mueller@jku.at), who hosted my research stay in 2026.

Thank you for considering my application.

Sincerely,\\
Grover Aruquipa

\end{document}
